\section{SIMULATION RESULTS}\label{sec:simulation-results}
To verify the effectiveness of the streaming federated learning framework for satellite system fault diagnosis, 
this paper designs a simulation experiment with 1 central server node and 6 TT\&C center nodes.
The central server is a logical entity responsible for coordinating the federated learning process, including distributing global models, aggregating local model updates, and maintaining partial models from communication-constrained TT\&C centers. It can be implemented either as a dedicated physical server or as a virtual node in cloud environments.
Among the TT\&C centers, 5 centers have no communication constraints, 
while 1 TT\&C center has communication constraints (this configuration aligns with the scale of China's TT\&C centers).

\added[id=cuiys]{The fault data used in this study is generated through a Simulink-based simulation framework~\cite{liu2015fault} that is calibrated against real telemetry from real satellite system.
Specifically, we first perform comparative analysis between simulated signals and measured telemetry, and then tune the simulation parameters to align key characteristics, 
so that the generated data preserves realistic fault evolution patterns.}
During the data preprocessing stage, 
data augmentation techniques are integrated to broaden the training data distribution 
\added[id=cuiys]{and improve model robustness, with the explicit goal of mimicking realistic on-orbit operating conditions 
and sensing disturbances. Four domain-specific approaches are developed for satellite systems:}
(1) Random noise addition to simulate space environment interference;
(2) Random temporal shifts to mimic timing offsets in anomaly detection;
(3) Amplitude scaling to simulate variations in lighting and temperature conditions;
(4) Temporal masking to simulate telemetry data loss during operation.
\added[id=cuiys]{These augmentation techniques are combined into a unified pipeline to emulate practical scenarios 
that are difficult to fully capture in collected data, 
thereby enabling the model to better adapt to diverse temporal patterns and complex space environments encountered in real satellite operations.}

To emulate the Non-IID characteristics of fault data across TT\&C centers, 
we employ the Dirichlet distribution for data partitioning, where parameter $\alpha$ controls the distribution sharpness~\cite{hsu2019measuring}. 
Fig.~\ref{fig:alpha_comparison} visualizes how the dataset partitioning outcomes change as $\alpha$ varies.

\begin{figure*}[t]
    \centering
	\includegraphics[width=\textwidth]{fig/Simulation Results/alpha_comparison.pdf}
    \caption{Dataset partitioning results under different $\alpha$ values.}
    \label{fig:alpha_comparison}
\end{figure*}

Throughout the entire training procedure, 
the Adam optimizer \cite{2014Adam} is employed with Dropout regularization \cite{srivastava2014dropout} 
applied to fully connected layers to enhance model generalization capability. 
Xavier initialization \cite{glorot2010understanding} is utilized for initializing weights between newly added output layer nodes 
and their preceding layers. 
All experimental evaluations are conducted on a computing cluster 
comprising 8 NVIDIA GeForce RTX 3080 Ti GPUs, 
with the framework implemented using the PyTorch deep learning library.

% \added{In standard federated learning (corresponding to Stage 1), 
% 60\% of TT\&C centers are randomly selected for each training round, 
% with each center conducting 5 epochs of local model training on their respective datasets using a learning rate of 0.001.}

% \added{For model training at the communication-constrained TT\&C center (corresponding to Stage 3), 
% the training configuration comprises 100 rounds with a local batch size of 16 samples, 
% and a reduced learning rate of 0.0001 to accommodate the constrained communication environment.}

% \added{In federated learning among unconstrained TT\&C centers (corresponding to Stage 5), 
% the training protocol spans 300 global rounds with 60\% center participation per round. 
% Each participating center executes 5 local training epochs per federated round using a learning rate of 0.0001, 
% which is deliberately reduced from the initial training phase to ensure stable convergence 
% and optimal performance within practical epoch constraints.}

% \added{All simulations in this paper are conducted under an extreme yet representative scenario: 
% the communication-constrained TT\&C center simultaneously hosts fault categories 
% that are entirely disjoint from those present in other TT\&C centers. 
% This configuration mimics challenging real-world conditions where certain ground stations encounter unique fault patterns 
% not observed elsewhere in the network.}

\subsection{Performance Evaluation}
Under this configuration, the proposed framework is deployed 
to perform fault diagnosis for satellite systems across all participating TT\&C centers. 
The experimental results are illustrated in Fig. \ref{fig:performance_across_datasets}. 
Three distinct datasets are constructed for comprehensive evaluation: 
Dataset 1 contains fault data exclusively from unconstrained TT\&C centers; 
Dataset 2 includes fault data solely from the communication-constrained TT\&C center; 
Dataset 3 represents the aggregated dataset comprising all fault instances from every TT\&C center 
(containing nine fault classes in total, with the ninth class being exclusive to the constrained TT\&C center).

Fig. \ref{fig:performance_across_datasets} illustrates the model's behavior across multiple datasets through different training phases. 
The outcomes highlight two principal capacities of the proposed framework:
First, during the learning phase at communication-constrained TT\&C centers, 
the dual-domain knowledge distillation-based incremental learning approach empowers the model to efficiently gain internal expertise 
while predominantly maintaining formerly acquired external knowledge (as demonstrated by consistent performance on Dataset 1). 
Second, during the learning phase at unconstrained TT\&C centers, 
the forgetting compensation functionality offered by the proxy server permits the model to continually integrate fresh external knowledge 
while efficiently preserving internal knowledge (as confirmed by persistent performance on Dataset 2).  
When benchmarked against ``joint training" with all fault data centralized at a single TT\&C center—
representing the theoretical performance upper bound—the proposed method achieves competitive results.
An important distinction to note is that the training timeline depicted in Fig.~\ref{fig:performance_across_datasets} 
should not be confused with federated learning rounds; specifically, the communication-constrained TT\&C center 
engages in only a single federated learning round throughout the entire process.

\begin{figure}[t]
	\centering
	\includegraphics[width=\columnwidth]{fig/Simulation Results/main result.pdf}
	\caption{Performance across datasets.}
	\label{fig:performance_across_datasets}
\end{figure}

In practical engineering deployments, satellites continuously produce novel fault patterns over their operational lifetime. 
To examine how the proposed framework behaves as new faults emerge over time, 
we adopt the streaming federated learning scenario schematized in Fig.~\ref{fig:streaming_scenario}. 
The experiment unfolds in two phases.
In the first phase, a new fault class appears: the unconstrained TT\&C centers jointly hold seven fault classes, 
whereas the communication-constrained TT\&C center contains one fault class that is completely disjoint from those at the other centers. 
In the second phase, an additional fault class emerges; the unconstrained centers retain the same seven fault classes, while the constrained center now hosts two fault classes that remain unique relative to the unconstrained centers. 
This configuration mimics the continuous arrival of novel fault types during real satellite operations.
\begin{figure}[t]
    \centering
    \includegraphics[width=\columnwidth]{fig/Simulation Results/streaming_scenario.pdf}
	\caption{Streaming federated learning scenario: fault class distribution across TT\&C centers.}
    \label{fig:streaming_scenario}
\end{figure}

Fig. \ref{fig:streaming_federated_learning_performance} illustrates the performance evolution of the algorithm 
under streaming federated learning conditions. 
The proposed framework maintains robust performance continuity 
when encountering novel fault classes at communication-constrained TT\&C centers 
or when incorporating newly constrained centers into the federated learning process. 
Critically, the framework sustains consistent performance levels across iterative streaming sessions, 
exhibiting minimal degradation over successive iterations given sufficient model capacity. 
These empirical findings establish essential theoretical groundwork and provide compelling evidence 
for the framework's practical deployability in real-world satellite systems.

\begin{figure}[t]
	\centering
	\includegraphics[width=\columnwidth]{fig/Simulation Results/Scalability.pdf}
	\caption{Streaming federated learning performance.}
	\label{fig:streaming_federated_learning_performance}
\end{figure}

\subsection{Comparative Experiments}
We conduct a comprehensive performance evaluation by comparing the proposed approach with 
established fault diagnosis methodologies under identical experimental conditions. 
The comparative analysis focuses on three critical metrics: classification accuracy, 
communication efficiency (measured in federated learning rounds), and robustness to data heterogeneity. 
Fig. \ref{fig:comparison_with_sota_satellite_fault_diagnosis} presents the performance comparison between our proposed algorithm 
and contemporary satellite system fault diagnosis approaches, 
particularly the state-of-the-art Dual-aspect Time Series Transformer (DTST) method~\cite{xu2024dtst}, 
which represents a specialized space system fault diagnosis technique.
To ensure equitable comparison conditions, we acknowledge that in practical deployments, 
each TT\&C center typically encounters only a subset of all possible fault patterns. 
However, for the purpose of rigorous benchmarking, we adopt a modified experimental setup where all fault types are assigned to a single TT\&C center 
while maintaining extremely sparse sample representation for certain fault classes. 
This configuration enables meaningful performance assessment while preserving the essential characteristics of real-world data distribution challenges.

\begin{figure}[t]
	\centering
	\includegraphics[width=\columnwidth]{fig/Simulation Results/compare-1.pdf}
	\caption{Comparison with SOTA satellite fault diagnosis.}
	\label{fig:comparison_with_sota_satellite_fault_diagnosis}
\end{figure}

As demonstrated, the DTST algorithm, when trained on limited fault datasets, 
exhibits significant challenges in learning from data-scarce fault samples, consequently resulting in constrained recognition capabilities. 
Although such performance may satisfy conventional evaluation metrics (as many studies require $\ge 80$\% accuracy for fault diagnosis), 
it proves inadequate for practical deployment due to complete recognition failure for certain critical fault classes,
as empirically validated by the confusion matrix in Fig.~\ref{fig:confusion-matrix-dtst}. 
In contrast, our proposed methodology successfully amalgamates fault knowledge from all participating TT\&C centers, 
systematically overcoming the inherent limitations imposed by single-center data constraints regarding both fault diversity and sample sufficiency,
as illustrated in Fig.~\ref{fig:confusion-matrix-proposed}.

\begin{figure}[t]
	\centering
	\includegraphics[width=0.8\columnwidth]{fig/Simulation Results/confusion-matrix-20250328.pdf}
	\caption{Confusion matrix (DTST method).}
	\label{fig:confusion-matrix-dtst}
\end{figure}
\begin{figure}[t]
	\centering
	\includegraphics[width=0.8\columnwidth]{fig/Simulation Results/confusion-matrix-4-effectiveness.pdf}
	\caption{Confusion matrix (Proposed method).}
	\label{fig:confusion-matrix-proposed}
\end{figure}

Fig.~\ref{fig:comparison_with_federated_learning_methods} presents the performance comparison results between the proposed algorithm and the base Fed-Avg algorithm~\cite{mcmahan2017communication}. 
Experimental results demonstrate that compared to the traditional multi-iteration Fed-Avg approach, 
the proposed algorithm requires only a single federated learning iteration with the communication-constrained TT\&C center to train a well-performing and stable model. 
In contrast, when the standard Fed-Avg algorithm operates with only a single iteration involving the communication-constrained TT\&C center, 
its performance degrades significantly after that center ceases participation in federated learning due to the lack of forgetting compensation mechanisms. 
Additionally, due to the extreme heterogeneity of the data, the Fed-Avg algorithm exhibits highly unstable performance during the training process. 
However, when the proposed method is incorporated into the multi-iteration Fed-Avg algorithm, the stability of the overall approach can be significantly enhanced. 
This demonstrates that the suggested approach can also enhance the robustness of federated learning algorithms effectively.

\begin{figure}[t]
	\centering
	\includegraphics[width=\columnwidth]{fig/Simulation Results/compare-2.pdf}
	\caption{Comparison with federated learning methods.}
	\label{fig:comparison_with_federated_learning_methods}
\end{figure}


\subsection{Ablation Studies}
To rigorously validate the effectiveness of the proposed framework, 
comprehensive ablation studies were systematically conducted on two core components:
the incremental learning algorithm based on dual-domain knowledge distillation 
and the federated learning algorithm based on partial forgetting compensation. 
Table \ref{tab:ablation-study} presents the final accuracy results 
of the incremental learning algorithm ablation study based on dual-domain knowledge distillation in Stage 3. 
The experimental results demonstrate that fine-tuning methodologies suffer from catastrophic forgetting, 
causing significant performance degradation.
Feature extraction techniques preserve historical knowledge but limit the capacity to learn novel information, 
resulting in suboptimal fault discrimination performance. 
In contrast, our proposed incremental learning approach effectively acquires new knowledge 
while preserving existing knowledge integrity, 
achieving superior classification accuracy across all fault categories.

\begin{table}[t]
\caption{Ablation study for federated learning components}
\centering
\label{tab:ablation-study}
\resizebox{\columnwidth}{!}{
\begin{tabular}{c|c|c}
\hline
 & \textbf{Method} & \textbf{Accuracy} \\ \hline
\multirow{3}{*}{Stage 3} 
 & IL-DDKD & \textbf{0.978} \\ \cline{2-3}
 & Fine-Tuning & 0.933\\ \cline{2-3}
 & Feature Extraction & 0.894\\ \hline
\multirow{3}{*}{Stage 5} 
 & Heritage + Dynamic weight & \textbf{0.974}\\ \cline{2-3}
 & Heritage & 0.898\\\cline{2-3}
 & Dynamic weight & 0.958\\ \hline
\end{tabular}
}
\end{table}

Table \ref{tab:ablation-study} also provides the final accuracy outcomes 
of the federated learning algorithm ablation study grounded in partial forgetting compensation in Stage 5. 
The experimental findings indicate that the proxy server architecture introduced in this paper 
attains better performance maintenance compared to individual components.

\subsection{Robustness Analysis}
This section assesses the robustness of the suggested algorithm 
under different configurations of TT\&C center numbers and data heterogeneity degrees, 
thereby providing comprehensive insights into the framework's adaptability across diverse operational environments.

\begin{figure}[t]
	\centering
	\includegraphics[width=\columnwidth]{fig/Simulation Results/robustness-1.pdf}
	\caption{Performance stability across different $\alpha$ values.}
	\label{fig:algorithm_performance_vs_heterogeneity}

\end{figure}

Fig. \ref{fig:algorithm_performance_vs_heterogeneity} illustrates the algorithm's performance characteristics 
under varying degrees of data heterogeneity, as controlled by the parameter $\alpha$. 
Empirical results demonstrate that the proposed algorithm maintains robust performance stability for $\alpha \ge 1$, 
with only marginal performance degradation observed at $\alpha = 0.1$. 

\begin{figure}[t]
	\centering
	\includegraphics[width=\columnwidth]{fig/Simulation Results/robustness-2.pdf}
	\caption{Performance stability across different TT\&C centers}
	\label{fig:algorithm_performance_vs_client_count}
\end{figure}

Fig. \ref{fig:algorithm_performance_vs_client_count} illustrates the algorithm's performance characteristics 
under varying numbers of TT\&C centers, thereby evaluating its scalability properties. 
Experimental results demonstrate that the proposed algorithm exhibits minimal performance variation 
across different client configurations, indicating robust insensitivity to client count fluctuations. 
This superior scalability characteristic constitutes a fundamental advantage for large-scale deployment scenarios 
and provides essential technical foundation for future engineering implementations in diverse satellite network architectures.